\documentclass[11pt]{article}

\usepackage[utf8]{inputenc}
\usepackage[T1]{fontenc}
\usepackage[a4paper,margin=2.5cm]{geometry}
\usepackage{amsmath,amssymb}
\usepackage{booktabs}
\usepackage{array}
\usepackage[hypertexnames=false]{hyperref}
\usepackage{float}

\title{Nested Economies: Resources, Transformations, and Concentration}
\author{Alan Seroul, Gary Klajer}
\date{April 2026}

\begin{document}

\maketitle

\begin{abstract}
We study a dynamic model of economic agents organized in hierarchical layers. Each agent has resources, transformation capacities, and preferences. Agents can exchange resources, capacities, and preferences with their direct neighbors in the hierarchy, either toward their subcomponents or toward their supersets. The goal is not to provide a complete realistic description of an economy, but to build a minimal laboratory for observing how local rules can produce concentration, specialization, instability, oligarchy, or redistribution. The main contribution is twofold: first, to formalize a hierarchical framework that can be simulated at large scale; second, to show experimentally that resource transformations become interpretable only in a conservative regime where they are actually active. In this regime, the stoichiometric mode is the most promising candidate for reducing wealth concentration.
\end{abstract}

\section{Motivation}

Many real systems are made of nested units. An individual belongs to a family, a company, a city, a region, and a country. A cell contains organelles, is embedded in cellular tissue, and then in an organism. An atom becomes part of molecules, which in turn enter into larger structures. In each of these examples, the relevant objects do not all live at the same scale.

Yet formal models often focus on a single scale of interaction. We describe individuals interacting with individuals, countries with countries, companies with companies, but it is rarer to describe interactions between levels within a single system. This omission is problematic. Major collective phenomena arise precisely from the mismatch between local objectives and the objectives of the structures that contain agents. A country may go to war even though most of its citizens do not want to die. An organization may maximize its survival or wealth while degrading the situation of some of its members. A subgroup may capture the resources of a larger whole without this concentration having been explicitly programmed.

To formalize this idea, we start from a set-based representation. An agent can be viewed as a set, and its components as subsets. A company may contain teams, which contain individuals. An organism may contain tissues, which contain cells. In this reading, the fundamental relation is not only an interaction relation, but an inclusion relation: an agent can be a component of a larger agent.

There is then a natural translation between systems of sets and directed graphs. Each set is represented by a node. An arrow $A \to B$ means that $A$ is included in $B$, or that $A$ contributes directly to $B$. To avoid self-membership paradoxes and loops of the form $A \in B \in A$, the corresponding graph must be acyclic. Directed acyclic graphs, or DAGs, therefore provide the most general version of this hierarchical representation. In this report, sets mainly serve to motivate the structure of the model; we do not yet use a full algebra of set operations.

The overall objective is to build a dynamic framework for these hierarchical agents. We do not aim to capture the full richness of social, biological, or physical systems with a single model. Rather, we want to isolate a few mechanisms that are simple enough to be simulated, but expressive enough to produce macroscopic effects.

The central question of the report is economic: under which local rules does a hierarchy of agents exchanging resources produce concentration, specialization, or redistribution? Are there modes of resource transformation that favor or limit oligarchy? What changes when only elementary agents are allowed to transform resources, rather than the whole hierarchy?

\section{Sets and Hierarchical Graphs}

The DAG representation is the most general one: it allows the same agent to belong to several larger agents. For example, an individual may belong to a company, a city, an association, and a family. Formally, a node may therefore have several parents.

Stronger restrictions can then be imposed.

\paragraph{Layered DAGs.}
A layered DAG imposes a notion of level. Nodes are partitioned into layers $0,1,\ldots,L-1$. An edge goes only from one level to the immediately higher level. This means that a microscopic agent cannot directly belong to a macroscopic agent without passing through intermediate-level agents. For example, an individual does not directly reach the level of a country in the representation: depending on the modeling choice, the path may pass through a family, a company, a city, or a region.

In this work, the level of an agent is therefore defined as the index of its layer. Layer $0$ contains the most elementary agents in the model. Layer $l+1$ contains agents that are more aggregated than those of layer $l$. An edge between layer $l$ and layer $l+1$ indicates that a node at level $l$ interacts with a more aggregated node at level $l+1$. Interactions are local: a node exchanges only with the immediately higher and immediately lower layers.

A natural extension would be to allow edges between non-adjacent layers, as long as the graph remains acyclic. This would introduce hierarchical skip connections: an agent could interact directly with a structure located two or three levels above it, without passing through every intermediate level. We do not keep this option in order to preserve a strict notion of scale neighborhood, but it provides an important degree of freedom for future extensions.

\paragraph{Trees.}
A tree is an even stronger restriction. Each node has at most one parent. This imposes a unitary inclusion relation: an agent can directly belong to only one larger agent. This assumption may be reasonable in some administrative settings, for instance if one requires an individual to have only one country of attachment, or an employee to work for only one company. It is, however, too restrictive if one wants to represent multiple memberships.

Here we use layered DAGs, which preserve the possibility of several parents while imposing scale proximity.

To enable large-scale simulations, we do not consider transitions that would explicitly modify inclusion relations. In other words, the graph structure is not rewritten over time: a node does not sequentially create or delete discrete memberships. This choice is important because structure-modification operations are strongly sequential and difficult to parallelize.

Instead, we fix a complete structure between adjacent layers: each node in a layer is connected to all nodes in the immediately higher layer, and to all nodes in the immediately lower layer when it exists. Membership is therefore not binary. We will speak of \textbf{weighted membership}. In a sense, a node belongs to all accessible higher-order structures, but to varying degrees. These degrees are carried by the learned allocation proportions. A weak inclusion relation corresponds to a proportion close to zero; a strong relation corresponds to a high proportion.

This interpretation is operational: between two adjacent levels, belonging to a structure means exchanging with it. If a cell carries out almost all of its exchanges with a given organism, then it belongs to that organism in the model. The cells of an organism exchange primarily with their host organism, not directly with other organisms. The model therefore replaces discrete membership with an exchange intensity between neighboring levels.

This approximation remains parallelizable, but it is not free. Between two adjacent layers, the cost of the flows is on the order of $O(N_l N_{l+1} n)$, where $n$ is the number of resources. Much larger simulations will probably require a sparse version of this weighted membership.

Layer $l$ contains $N_l$ nodes.

Two topologies are used in the experiments:
\begin{itemize}
    \item \textbf{Flat topology}: $[10,10,10,10,10,10,10,10,10,10]$. All layers have the same size. The term flat means here that the width does not decrease with height; the structure is still a layered hierarchy.
    \item \textbf{Hier topology}: $[32,16,8,4,2,1]$. The number of nodes decreases with height, up to a single top node.
\end{itemize}

\section{Economic Model}

Economic theory provides a natural basis for this project, because it already works with agents, resources, preferences, and allocation decisions. Even when agents are not individuals, but companies, households, institutions, or states, economics provides a language for describing their constraints, objectives, and exchanges.

In neoclassical economics, three concepts play a central role: capital, labor, and utility. Capital denotes the resources or assets available for production, exchange, or consumption. Labor denotes an agent's capacity for action or production. Utility describes an agent's own satisfaction, and therefore what motivates its choices. Our model draws directly on these three notions. For each node, we represent capital by a resource vector, labor by a skill matrix, and utility by a preference vector.

The choices that follow nevertheless remain modeling choices. They are often partly arbitrary, and their relevance depends strongly on the scale considered. A rule that is natural for representing companies may be poor for representing cells; a rule that is reasonable for individuals may become absurd for states. To limit conclusions that would depend too much on a particular convention, we have therefore implemented several versions of natural mechanisms. The objective is to separate what is specific to a given mechanism from what seems structural in the hierarchical system itself.

\subsection{Node State}

Each node $i$ in layer $l$ has three objects:
\begin{align}
    r_i^{(l)}(t) &\in \mathbb{R}_+^{n}, \\
    S_i^{(l)}(t) &\in \mathbb{R}_+^{n \times n}, \\
    u_i^{(l)}(t) &\in \Delta^{n-1}.
\end{align}

Here $n$ is the number of resource types. The vector $r_i^{(l)}$ contains the quantities of resources held by the node. The matrix $S_i^{(l)}$ encodes its transformation capacities. The vector $u_i^{(l)}$ encodes its preferences, with $\sum_m u_{i,m}^{(l)} = 1$.

The economic interpretation is direct:
\begin{itemize}
    \item resources play the role of capital or available stocks;
    \item skills represent conversion, production, or organizational capacities;
    \item preferences describe what makes an agent satisfied.
\end{itemize}

For example, in a very simplified economy, resources could be food, energy, and tools. Skills would then describe an agent's capacity to transform energy and tools into food, or to convert certain resources into others. Preferences would indicate whether the agent values food, energy, or tools more. The same formalism can then be read at another scale: a company, a city, or an organism also has stocks, transformation capacities, and its own objectives.

\subsection{Flows Between Nodes}

At each step, a node can send a fraction of its resources, skills, and preferences to its parents or to its children. We denote by
\begin{align}
    P^{R,\uparrow,l}_{i,j,m}, \quad P^{R,\downarrow,l}_{i,k,m}
\end{align}
the proportions of resource $m$ sent by node $i$ in layer $l$ to a parent $j$ in layer $l+1$, or to a child $k$ in layer $l-1$. The same objects exist for skills, denoted $P^S$, and for preferences, denoted $P^U$.

The proportions are produced by optimizable logits. By construction, they are constrained to remain in $[0,1]$ and to respect an allocation budget.

\subsection{Budget Modes}

Two budget modes are implemented.

\paragraph{Joint budget.}
For each node and each resource, a single softmax allocates the budget between three groups: keep, send upward, send downward. Formally, if the available groups are $G \subseteq \{\mathrm{keep},\uparrow,\downarrow\}$, the group budget is:
\begin{equation}
    b_g = \frac{\exp a_g}{\sum_{h \in G} \exp a_h}.
\end{equation}
Within the $\uparrow$ or $\downarrow$ group, a second softmax distributes the mass among the available partners.

\paragraph{Split budget.}
When the node has both parents and children, the upward budget and the downward budget are decided separately. A keep/up softmax is applied for parents and a keep/down softmax for children. This mode therefore allows stronger simultaneous pressure in both directions.
It should not be read as a single physical budget that is strictly conserved at the node level, but as a more permissive allocation rule between upward and downward flows.

\subsection{Net Flows}

The resource flows of a node decompose into an intrinsic part, linked to children, and an extrinsic part, linked to parents:
\begin{align}
    r_{\mathrm{int},i}^{(l)} &=
    \sum_k P^{R,\uparrow,l-1}_{k,i} \odot r_k^{(l-1)}
    -
    \sum_k P^{R,\downarrow,l}_{i,k} \odot r_i^{(l)}, \\
    r_{\mathrm{ext},i}^{(l)} &=
    \sum_j P^{R,\downarrow,l+1}_{j,i} \odot r_j^{(l+1)}
    -
    \sum_j P^{R,\uparrow,l}_{i,j} \odot r_i^{(l)}.
\end{align}

We set:
\begin{equation}
    r_{\mathrm{net},i}^{(l)} = r_{\mathrm{int},i}^{(l)} + r_{\mathrm{ext},i}^{(l)}.
\end{equation}
Skills and preferences follow the same principle. For skills, the proportions act on the rows of $S$, that is, on the capacities associated with each input resource.

This choice also gives an important intuition about collaboration. Combining two agents is not merely adding two separate productions. If two agents pool resources $r_1,r_2$ and skills $S_1,S_2$, then in general
\begin{equation}
    (r_1+r_2)(S_1+S_2) \neq r_1S_1+r_2S_2.
\end{equation}
The combination of capacities can therefore create nonlinear effects, positive or negative. This remark does not modify the simulation rule, but it explains why it is relevant to circulate skills in addition to resources.

\section{Model Modes}

\subsection{Resource Modes}

Three resource modes are important.

\paragraph{\texttt{zero\_sum\_no\_s}.}
In this mode, there is no internal transformation of resources. The matrix $S$ is frozen and ignored in the dynamics of $r$. The update is:
\begin{equation}
    r_i(t+1) = \operatorname{ReLU}\left(r_i(t) + r_{\mathrm{net},i}(t)\right).
\end{equation}
This mode corresponds to pure exchanges. It conserves mass through redistribution, except for positive-projection effects when a component would become negative. It served as the initial battery. It is useful for studying the effects of topology, budget, number of resources, and optimization mode.

\paragraph{\texttt{zero\_sum\_stoch\_s}.}
In this mode, resources are transformed by $S$, then redistributed. After the update, the matrices $S$ are normalized row by row. In addition, the total resource mass is globally rescaled to conserve the total mass of the system. This mode activates transformations while avoiding unbounded growth:
\begin{equation}
    r_i(t+1) = \operatorname{ReLU}\left(T(r_i(t),S_i(t)) + r_{\mathrm{net},i}(t)\right).
\end{equation}
This is the cleanest regime for comparing transformation modes. Conservation is global: it concerns the total mass of the system after rescaling, not the wealth of each node, nor necessarily each resource type taken separately.

\paragraph{\texttt{free}.}
In this mode, transformations are active, but no global mass conservation is imposed. Skills can evolve freely and are not normalized as in the stochastic mode. This mode represents a non-conservative productive or destructive regime. It is informative for observing growth, extinction, and instability, but it should not be compared directly with conservative modes without an appropriate metric.

\subsection{Transformation Modes}

We denote the internal resource transformation by $T(r,S)$.

\paragraph{\texttt{markov}.}
The transformation is linear:
\begin{equation}
    T_m(r,S) = \sum_n r_n S_{n,m}.
\end{equation}
When $S$ is row-stochastic, this transformation redistributes mass among resource types.

\paragraph{\texttt{mass\_action}.}
The transformation is multiplicative in log-space:
\begin{equation}
    \log T_m(r,S) = \sum_n S_{n,m} \log(r_n + \varepsilon).
\end{equation}
Thus:
\begin{equation}
    T_m(r,S) = \exp\left(\sum_n S_{n,m}\log(r_n+\varepsilon)\right).
\end{equation}
This mode resembles multiplicative kinetics: several resources can jointly contribute to producing a resource type.

\paragraph{\texttt{stoichiometric}.}
We keep the English name used in the code. Here it denotes reaction dynamics with stoichiometric coefficients.

We introduce $K$ reactions. Each reaction $k$ has an input vector $W^{\mathrm{in}}_k$ and an output vector $W^{\mathrm{out}}_k$, both positive. For an agent $i$, the raw reaction rate is $v_{i,k}$. The effective rate is limited by the limiting reactant:
\begin{equation}
    v^{\mathrm{max}}_{i,k} = \min_n \frac{r_{i,n}}{W^{\mathrm{in}}_{k,n}+\varepsilon},
    \qquad
    v^{\mathrm{actual}}_{i,k} = \min(v_{i,k},v^{\mathrm{max}}_{i,k}).
\end{equation}
The resource variation is:
\begin{equation}
    \Delta r_i = \sum_k v^{\mathrm{actual}}_{i,k}
    \left(W^{\mathrm{out}}_k - W^{\mathrm{in}}_k\right).
\end{equation}
The transformation returns $T(r_i)=r_i+\Delta r_i$.

\paragraph{\texttt{leontief}.}
Each produced resource $m$ has a recipe of prerequisites $C_{j,m}$. The production scale is:
\begin{equation}
    q_{i,m} = \min_j \frac{r_{i,j}}{C_{j,m}+\varepsilon}.
\end{equation}
Effective production is modulated by an efficiency $e_{i,m}\in[0,1]$:
\begin{equation}
    T_m(r_i)=q_{i,m} e_{i,m}.
\end{equation}
This mode imposes strong complementarity among resources. A missing resource can block an entire production process.

\subsection{\texttt{ATOMS\_ONLY\_TRANSFORM}}

This parameter decides where the internal transformation $T$ is applied.
\begin{itemize}
    \item If \texttt{ATOMS\_ONLY\_TRANSFORM=True}, only nodes in layer $0$ transform resources. Higher layers redistribute without transforming.
    \item If \texttt{ATOMS\_ONLY\_TRANSFORM=False}, all nodes apply their internal transformation.
\end{itemize}
This parameter has no effect in \texttt{zero\_sum\_no\_s}, since the transformation is bypassed. It becomes decisive in the other two modes.

\subsection{Optimization Modes}

\paragraph{\texttt{reciprocity}.}
Each agent keeps an exponential memory of useful resources received from its partners. If $M_{i,j,m}$ is the memory of what $i$ received from $j$ for resource $m$, then:
\begin{equation}
    M_{i,j,m}(t+1)=\alpha M_{i,j,m}(t)+(1-\alpha)u_{i,m}(t) \cdot \mathrm{received}_{i,j,m}(t).
\end{equation}
The optimizer receives a bonus when an agent sends resources to partners that have historically sent useful resources back to it. This mechanism avoids a strictly opportunistic strategy in which each agent would keep everything in the short term. The intuition is that, without memory, an agent may be incentivized not to remunerate its collaborators at the current step, even if this strategy later destroys the structure that enabled it to produce. Memory therefore forces the model to distinguish an immediate gain from a durable relation.

\paragraph{\texttt{market}.}
In this mode, resource routing is biased toward partners whose resources or preferences are aligned with the sending agent. Upward, the bonus depends on $u_i \odot r_j$: an agent prefers to send to a parent that has the resources it values. Downward, the bonus depends on $r_i \odot u_k$: an agent prefers to send to a child that values what it has.

These two modes are differentiable objective terms. They give a local behavioral interpretation to the flows, but the simulation remains synchronously optimized by a global optimizer. It is therefore not yet a decentralized game in which each agent independently solves its own problem.

\section{Utility and Optimization}

An agent's utility is a Cobb-Douglas function:
\begin{equation}
    U_i^{(l)} = \prod_m \left(r_{i,m}^{(l)}\right)^{u_{i,m}^{(l)}}.
\end{equation}
In practice, we optimize log-utility:
\begin{equation}
    \log U_i^{(l)} = \sum_m u_{i,m}^{(l)}\log(r_{i,m}^{(l)}+\varepsilon).
\end{equation}

The aggregate objective is the sum of mean utilities by layer:
\begin{equation}
    \mathcal{L}_{\mathrm{primary}} = -\sum_l \frac{1}{N_l}\sum_i \log U_i^{(l)}.
\end{equation}
The negative sign comes from the fact that the optimizer minimizes the loss. In \texttt{reciprocity} mode, a reciprocity term is added. In \texttt{market} mode, a market-alignment term is added. The optimized parameters are the logits that determine the allocation proportions.

This optimization choice is pragmatic. One could seek an exact allocation rule, but the problem combines several resources, several agents, upward and downward flows, node-specific preferences, and temporal dynamics. It therefore does not correspond directly to the classical static allocation problems for which standard analytical solutions are available. Gradient descent provides a uniform method for adjusting interaction proportions in this differentiable framework. It does not claim to be a complete description of individual decision-making; here it serves as a search mechanism for local policies compatible with the aggregate objective.

\section{Metrics}

The results are evaluated with the following metrics.

\paragraph{Final Gini.}
The Gini coefficient is computed on each node's total wealth:
\begin{equation}
    W_i=\sum_m r_{i,m}.
\end{equation}
It measures wealth concentration.

\paragraph{Top20 share.}
Share of total wealth held by the richest $20\%$ of nodes.

\paragraph{Oligarchy.}
In this report, we call oligarchy a situation in which wealth is highly concentrated in a small group of nodes. It is measured mainly by a high Gini and a high top20 share. When the same nodes remain in the top over time, we also speak of elite persistence.

\paragraph{Elite persistence.}
Persistence compares the set of nodes belonging to the top $20\%$ over time. Initial persistence measures the survival of the initial elite; local persistence measures the stability of the top between two consecutive snapshots.

\paragraph{Portfolio entropy.}
For each node, we normalize its resource portfolio:
\begin{equation}
    p_{i,m}=\frac{r_{i,m}}{\sum_n r_{i,n}+\varepsilon}.
\end{equation}
The mean entropy is:
\begin{equation}
    H=-\frac{1}{N}\sum_i\sum_m p_{i,m}\log(p_{i,m}+\varepsilon).
\end{equation}
It measures the average diversity of resources held by agents.

\paragraph{Final utility.}
Final average of log-utilities.

\paragraph{Numerical stability.}
In non-conservative runs, we also track the presence of non-finite values in the saved history and the fraction of finite values.

\section{Experimental Design}

The first battery contains $80$ runs:
\begin{equation}
    2 \text{ topologies} \times 8 \text{ variants} \times 5 \text{ seeds}.
\end{equation}
All these runs use \texttt{RESOURCE\_MODE=zero\_sum\_no\_s}. This battery makes it possible to study the effects of topology, budget, market mode, number of resources, and reciprocity.

A second battery contains $32$ runs:
\begin{equation}
    2 \text{ resource modes} \times 4 \text{ transformations}
    \times 2 \text{ atoms-only values} \times 2 \text{ topologies}.
\end{equation}
The tested modes are \texttt{zero\_sum\_stoch\_s} and \texttt{free}. Each combination is launched once. In this second battery, the behavioral mode is fixed to \texttt{reciprocity}, the budget to \texttt{joint}, and the number of resources to $5$. This second battery should therefore be read as a diagnostic exploration, not as a robust statistical validation.

\begin{table}[H]
\centering
\scriptsize
\begin{tabular}{lll}
\toprule
Factor & Values & Role \\
\midrule
RESOURCE\_MODE & zero\_sum\_no\_s, zero\_sum\_stoch\_s, free & conservation and activation of $S$ \\
TRANSFORM\_MODE & mass action, markov, stoichiometric, leontief & internal transformation mechanism \\
ATOMS\_ONLY & True, False & localization of transformation \\
TOPOLOGY & flat, hier & multi-scale structure \\
MODE & reciprocity, market & behavioral objective \\
BUDGET\_MODE & joint, split & allocation constraint \\
n\_res & 5, 15 & resource diversity \\
\bottomrule
\end{tabular}
\caption{Experimental factors manipulated in the two batteries.}
\end{table}

\section{Main Results}

\subsection{The Initial Battery Shows a Non-Effect of Transformations}

In \texttt{zero\_sum\_no\_s}, the configurations \texttt{markov}, \texttt{stoichiometric}, \texttt{leontief}, \texttt{mass\_action}, and \texttt{ATOMS\_ONLY\_TRANSFORM=False} produce the same metrics as the corresponding reference. This result is not an economic phenomenon, but a direct consequence of the mode definition: the internal transformation is disabled.

The initial battery remains informative on the other dimensions:
\begin{itemize}
    \item the hierarchical topology increases average concentration, with a mean Gini of $0.848$ compared with $0.728$ in the flat topology;
    \item increasing the number of resources to $15$ strongly reduces concentration in the flat topology, with a mean final Gini of $0.430$;
    \item the \texttt{market} mode produces the strongest concentrations, with a mean final Gini of $0.906$ in the hierarchical topology;
    \item the \texttt{split} mode increases flow exclusivity;
    \item elites are very stable locally, with local persistence close to $1$ in the reference configurations.
\end{itemize}

\subsection{\texttt{zero\_sum\_stoch\_s} Makes Transformations Comparable}

The central result comes from the complementary battery. In \texttt{zero\_sum\_stoch\_s}, all transformations remain numerically stable and total mass is conserved. The transformations are then no longer equivalent.

\begin{table}[H]
\centering
\begin{tabular}{lrrrr}
\toprule
Transformation & Gini & Top20 & Utility & Entropy \\
\midrule
mass action & 0.909 & 0.963 & -6.975 & 1.147 \\
markov & 0.910 & 0.967 & -6.769 & 1.187 \\
stoichiometric & 0.562 & 0.649 & 0.138 & 0.205 \\
leontief & 0.852 & 0.887 & -4.909 & 0.077 \\
\bottomrule
\end{tabular}
\caption{Final averages in \texttt{zero\_sum\_stoch\_s}.}
\end{table}

The reasoning is therefore as follows. First, the initial battery shows that \texttt{zero\_sum\_no\_s} cannot test transformations. Next, \texttt{zero\_sum\_stoch\_s} activates these transformations while conserving a comparable mass. Finally, in this clean regime, \texttt{mass\_action} and \texttt{markov} remain close: they produce high concentration and a large share of wealth held by the top $20\%$. \texttt{leontief} also maintains high concentration, with highly specialized portfolios. By contrast, \texttt{stoichiometric} strongly reduces concentration and gives the best final utility. This result does not mean that portfolios become more diversified: its mean entropy is low. The signal therefore concerns mainly the distribution of wealth among agents, not the internal diversity of the resources held by each agent.

\subsection{The \texttt{free} Mode Is a Growth or Extinction Regime}

In \texttt{free}, transformations are active but mass is not conserved. The resulting behavior is not directly comparable to the conservative regime. Some configurations produce enormous masses, on the order of $10^{36}$. Others die out completely.

The \texttt{free} mode should therefore be interpreted as a separate regime. In its current state, it cannot be used to draw quantitative conclusions about distribution. It mainly serves as a diagnostic: without conservation, cost, or capacity constraints, the dynamics can be dominated by numerical growth or extinction. A usable regime will have to add normalization or a capacity constraint.

\subsection{Effect of \texttt{ATOMS\_ONLY\_TRANSFORM}}

In \texttt{zero\_sum\_no\_s}, \texttt{ATOMS\_ONLY\_TRANSFORM} changes nothing. In modes where transformations are active, it becomes important. In \texttt{zero\_sum\_stoch\_s}, allowing all layers to transform reduces average concentration:
\begin{equation}
    Gini(\mathrm{atoms\ only}) = 0.885,
    \qquad
    Gini(\mathrm{all\ nodes}) = 0.731.
\end{equation}
In \texttt{free}, allowing all layers to transform often makes the dynamics more unstable, especially for \texttt{mass\_action}, \texttt{markov}, and \texttt{leontief}.

\section{Discussion and Limitations}

The proposed model provides a differentiable framework for studying nested economic agents. Its central point is to make agents coexist at several levels, with upward and downward flows of resources, skills, and preferences. Discrete inclusion is replaced by weighted membership, which makes the simulation parallelizable while preserving a hierarchical interpretation.

The results suggest that hierarchy is already sufficient to produce concentration, but that the internal transformation mechanism becomes decisive as soon as it is actually activated in a conservative regime. The most interesting signal concerns \texttt{stoichiometric}, which clearly reduces concentration in the complementary battery.

The model also shows that hierarchical aggregates are not merely larger individuals. A hierarchy can concentrate wealth at the top, but also create intermediate bottlenecks. The \texttt{market} mode finally indicates that local supply-demand alignment can reinforce already attractive nodes instead of redistributing resources.

Limitations.\\
The initial battery contains $5$ seeds per configuration, but on small graphs. The complementary battery contains only $1$ seed per combination. The results on the stochastic conservative mode, especially the interest of \texttt{stoichiometric}, must therefore be replicated.\\
The graph does not change over time: membership degrees evolve through allocation proportions, but the possible edges remain fixed between adjacent layers.\\
All parameters are optimized by Adam in a synchronous simulation. This is not yet a model of asynchronous agents each making independent decisions with their own time scale.\\
The \texttt{free} mode confirms a more general limitation: without capacity, cost, or conservation constraints, the results can be dominated by numerical growth or extinction. The main conclusions therefore rest on \texttt{zero\_sum\_stoch\_s}, not on \texttt{free}.

The runs can be explored qualitatively with the dashboard \texttt{viz.py}, for example with \texttt{python viz.py} or \texttt{python viz.py runs/<run\_id>.pkl}. The next experiments should mainly replicate \texttt{zero\_sum\_stoch\_s} with several seeds, keep \texttt{ATOMS\_ONLY\_TRANSFORM} as a factor, and set \texttt{free} aside until a capacity or normalization constraint is defined.

The contribution of the report is therefore twofold: a hierarchical economic model that can be simulated at large scale, and a first experimental separation between structure effects, conservation effects, and transformation effects.

\appendix

\section{Numerical Values of the Initial Battery}

\begin{table}[H]
\centering
\scriptsize
\begin{tabular}{llrrrr}
\toprule
Family & Topology & Gini & Top20 & Entropy & Corr U-W \\
\midrule
REF & flat & 0.750 & 0.821 & 0.447 & 0.850 \\
REF & hier & 0.851 & 0.922 & 0.317 & 0.955 \\
nres15 & flat & 0.430 & 0.519 & 0.709 & 0.331 \\
nres15 & hier & 0.743 & 0.787 & 0.439 & 0.678 \\
market & flat & 0.877 & 0.975 & 0.110 & 0.999 \\
market & hier & 0.906 & 0.997 & 0.081 & 0.999 \\
split & flat & 0.762 & 0.882 & 0.362 & 0.963 \\
split & hier & 0.875 & 0.969 & 0.176 & 0.998 \\
\bottomrule
\end{tabular}
\caption{Non-redundant summary of the initial battery. The markov, stoichiometric, leontief, and allnodes variants are not listed because they are identical to REF in \texttt{zero\_sum\_no\_s}. Values are final averages over 5 seeds.}
\end{table}

\section{Numerical Values of Active Transformations}

\begin{table}[H]
\centering
\scriptsize
\begin{tabular}{llrrrr}
\toprule
Resource mode & Transformation & Gini & Top20 & Utility & Entropy \\
\midrule
zero\_sum\_stoch\_s & mass action & 0.909 & 0.963 & -6.975 & 1.147 \\
zero\_sum\_stoch\_s & markov & 0.910 & 0.967 & -6.769 & 1.187 \\
zero\_sum\_stoch\_s & stoichiometric & 0.562 & 0.649 & 0.138 & 0.205 \\
zero\_sum\_stoch\_s & leontief & 0.852 & 0.887 & -4.909 & 0.077 \\
free & mass action & 0.713 & 0.727 & 30.903 & 1.035 \\
free & markov & 0.842 & 0.904 & 35.527 & 0.948 \\
free & stoichiometric & 0.824 & 0.846 & 8.381 & 0.772 \\
free & leontief & 0.322 & 0.332 & 35.990 & 0.523 \\
\bottomrule
\end{tabular}
\caption{Summary of the 32 complementary runs, aggregated by resource mode and transformation. Each average groups the two topologies and the two values of \texttt{ATOMS\_ONLY\_TRANSFORM}, with a single seed per combination. The \texttt{free} values are diagnostic, not clean estimates of an unconstrained free regime, because of numerical protections, explosions, and extinctions.}
\end{table}

\section{Effect of \texttt{ATOMS\_ONLY\_TRANSFORM}}

\begin{table}[H]
\centering
\scriptsize
\begin{tabular}{llrrrr}
\toprule
Resource mode & ATOMS\_ONLY & Gini & Final wealth & Utility & Entropy \\
\midrule
zero\_sum\_stoch\_s & True & 0.885 & 402.278 & -4.464 & 0.442 \\
zero\_sum\_stoch\_s & False & 0.731 & 402.278 & -4.794 & 0.866 \\
free & True & 0.845 & $4.34 \cdot 10^{16}$ & 22.529 & 0.817 \\
free & False & 0.532 & $2.04 \cdot 10^{36}$ & 32.872 & 0.846 \\
\bottomrule
\end{tabular}
\caption{Average effect of \texttt{ATOMS\_ONLY\_TRANSFORM}.}
\end{table}

\section{Unstable or Degenerate \texttt{free} Runs}

\begin{table}[H]
\centering
\scriptsize
\begin{tabular}{lll}
\toprule
Readable name & Configuration & Symptom \\
\midrule
free-mass-allnodes-flat & mass action, atoms=False, flat & final wealth $6.79 \cdot 10^{36}$ \\
free-markov-allnodes-flat & markov, atoms=False, flat & infinite values in the history \\
free-leontief-atoms-flat & leontief, atoms=True, flat & infinite values in the history \\
free-leontief-allnodes-flat & leontief, atoms=False, flat & total extinction \\
free-mass-allnodes-hier & mass action, atoms=False, hier & final wealth $4.16 \cdot 10^{36}$ \\
free-markov-allnodes-hier & markov, atoms=False, hier & final wealth $3.34 \cdot 10^{36}$ \\
free-leontief-atoms-hier & leontief, atoms=True, hier & final wealth $3.04 \cdot 10^{17}$ \\
free-leontief-allnodes-hier & leontief, atoms=False, hier & total extinction \\
\bottomrule
\end{tabular}
\caption{Problematic configurations of the \texttt{free} mode.}
\end{table}

\end{document}
